\documentclass[12pt,a4paper]{article}

% --- PACOTES BÁSICOS E CONFIGURAÇÃO ---
\usepackage[utf8]{inputenc}
\usepackage[T1]{fontenc}
\usepackage[brazil]{babel}
\usepackage[top=3cm, bottom=2cm, left=3cm, right=2cm]{geometry}
\usepackage{graphicx}
\usepackage{amsmath,amssymb}
\usepackage{booktabs}
\usepackage{float}
\usepackage{caption}
\usepackage{xcolor}
\usepackage{listings}
\usepackage{hyperref}
\usepackage{indentfirst}

% --- CONFIGURAÇÃO DE LINKS HIPERTEXTO ---
\hypersetup{
    colorlinks=true,
    linkcolor=blue!70!black,
    citecolor=blue!70!black,
    urlcolor=blue!70!black
}

% --- CONFIGURAÇÃO DO LISTINGS PARA CÓDIGO PYTHON ---
\definecolor{codegreen}{rgb}{0.13,0.54,0.13}
\definecolor{codegray}{rgb}{0.5,0.5,0.5}
\definecolor{codepurple}{rgb}{0.58,0,0.82}
\definecolor{backcolour}{rgb}{0.97,0.97,0.97}

\lstset{
    language=Python,
    backgroundcolor=\color{backcolour},
    commentstyle=\color{codegreen}\itshape,
    keywordstyle=\color{blue}\bfseries,
    numberstyle=\tiny\color{codegray},
    stringstyle=\color{codepurple},
    basicstyle=\ttfamily\footnotesize,
    breakatwhitespace=false,
    breaklines=true,
    captionpos=b,
    keepspaces=true,
    numbers=left,
    numbersep=7pt,
    showspaces=false,
    showstringspaces=false,
    showtabs=false,
    tabsize=4,
    frame=single,
    rulecolor=\color{black!20}
}

\begin{document}

% ==============================================================================
% CAPA DO RELATÓRIO
% ==============================================================================
\begin{titlepage}
    \centering
    {\large \textbf{UNIVERSIDADE FEDERAL DE VIÇOSA}}\\[0.2cm]
    {\large \textbf{CAMPUS FLORESTAL}}\\[0.2cm]
    {\large INSTITUTO DE CIÊNCIAS EXATAS E TECNOLÓGICAS}\\[0.2cm]
    {\large CURSO DE CIÊNCIA DA COMPUTAÇÃO}\\[3.5cm]

    {\bfseries \Large TRABALHO PRÁTICO I:}\\[0.3cm]
    {\bfseries \LARGE BIBLIOTECA DE ANÁLISE DE REDES SOCIAIS EM GRAFOS COM MATRIZ DE ADJACÊNCIA}\\[3.5cm]

    \begin{flushleft}
        \textbf{Disciplina:} CCF-331 -- Teoria e Modelo de Grafos\\
        \textbf{Professor:} Marcus Henrique Soares Mendes\\[0.8cm]
        \textbf{Equipe de Desenvolvimento:}\\
        Pedro Mendes Gonzaga -- Matrícula: 6563\\
        Otávio Soares Pedroso -- Matrícula: 6593\\
        Gabriel Marcos de Oliveira Félix -- Matrícula: [Inserir Matrícula]
    \end{flushleft}

    \vfill
    {\large Florestal -- MG}\\
    {\large Outubro de 2026}
\end{titlepage}

% ==============================================================================
% SUMÁRIO AUTOMÁTICO
% ==============================================================================
\tableofcontents
\newpage

% ==============================================================================
% 1. INTRODUÇÃO E MODELAGEM DO PROBLEMA
% ==============================================================================
\section{Introdução e Modelagem do Problema}

O avanço das plataformas digitais de relacionamento impulsionou a necessidade de ferramentas especializadas na análise estrutural e comportamental de redes de usuários. No escopo deste trabalho prático, foi projetada e implementada uma biblioteca modular em Python para manipulação, processamento e análise de redes sociais representadas formalmente por \textbf{grafos não direcionados e ponderados} $G = (V, E, w)$.

A modelagem matemática adotada estabelece que:
\begin{itemize}
    \item Cada \textbf{vértice} $v \in V$ representa um usuário da rede social.
    \item Cada \textbf{aresta} $\{u, v\} \in E$ representa uma conexão ou relacionamento direto mútuo entre os usuários $u$ e $v$.
    \item O \textbf{peso} $w(u, v) \in \mathbb{R}^+$ representa o grau de proximidade ou afinidade entre os usuários. Sob a perspectiva de otimização, quanto menor o peso numérico, menor é o ``custo'' de comunicação ou travessia entre eles, viabilizando o uso direto de algoritmos de caminho mínimo.
\end{itemize}

O objetivo primordial do projeto consiste no desenvolvimento de uma arquitetura estritamente modular, desacoplada e independente do programa principal, fundamentada obrigatoriamente na representação de \textbf{Matriz de Adjacência / Matriz de Valores}.

% ==============================================================================
% 2. ARQUITETURA DA BIBLIOTECA E ESTRUTURA DE DADOS
% ==============================================================================
\section{Arquitetura do Software e Estrutura de Dados}

A biblioteca foi decomposta em módulos especializados para assegurar coesão funcional, manutenibilidade e facilidade de teste:
\begin{enumerate}
    \item \textbf{\texttt{grafo.py}}: Contém a classe principal \texttt{Grafo}, o carregador de arquivos de texto e operações elementares.
    \item \textbf{\texttt{conectividade.py}}: Contém os algoritmos de busca em largura, cálculo de componentes conexas via Fechamento Transitivo de Roy, identificação de vértices de articulação e detecção de ciclos.
    \item \textbf{\texttt{caminhos\_analise.py}}: Contém os algoritmos de caminhos mínimos e rotas via Dijkstra.
    \item \textbf{\texttt{main.py}}: Interface de linha de comando com menu interativo interligando todos os módulos.
\end{enumerate}

\subsection{Representação por Matriz de Valores}
A estrutura interna foi alocada utilizando uma matriz bidimensional de dimensões $(N+1) \times (N+1)$, onde $N = |V|$ representa o número total de vértices. A linha e coluna de índice 0 foram propositalmente mantidas vazias para permitir a \textbf{indexação direta de 1 a $N$}, eliminando deslocamentos artificiais de índice (\textit{zero-based indexing}) e prevenindo falhas de limites.

A ausência de aresta entre dois vértices é representada pelo valor sentinela \texttt{None}, permitindo diferenciar com exatidão a inexistência de conexão de uma aresta com peso nulo. A listagem \ref{cod:grafo_init} ilustra a instanciação da estrutura e a adição simétrica de arestas não direcionadas.

\begin{lstlisting}[caption={Alocação da Matriz de Valores em \texttt{grafo.py}}, label={cod:grafo_init}]
class Grafo:
    def __init__(self, num_vertices: int):
        self.num_vertices = num_vertices
        # Matriz (N+1) x (N+1) inicializada com None
        self.matriz = [[None] * (num_vertices + 1) for _ in range(num_vertices + 1)]

    def adicionar_aresta(self, u: int, v: int, peso: float) -> None:
        # Grafo nao direcionado: atribuicao bidirecional simetrica
        self.matriz[u][v] = peso
        self.matriz[v][u] = peso
\end{lstlisting}

% ==============================================================================
% 3. OPERAÇÕES BÁSICAS E MÉTRICAS ESTRUTURAIS
% ==============================================================================
\section{Métricas Estruturais Básicas}

O módulo central disponibiliza as funções para extração de parâmetros fundamentais da topologia da rede:
\begin{itemize}
    \item \textbf{Ordem do Grafo ($|V|$):} Número total de usuários cadastrados na rede.
    \item \textbf{Tamanho do Grafo ($|E|$):} Quantidade de conexões únicas. Para evitar contagem duplicada na matriz simétrica, a biblioteca itera apenas pelo triângulo superior da matriz ($j \ge i$):
    \begin{equation}
        |E| = \sum_{i=1}^{N} \sum_{j=i}^{N} [M_{ij} \ne \text{None}]
    \end{equation}
    \item \textbf{Densidade $\varepsilon(G)$:} Razão entre o número de arestas e o número de vértices:
    \begin{equation}
        \varepsilon(G) = \frac{|E|}{|V|}
    \end{equation}
    \item \textbf{Densidade Relativa Normalizada ($D$):} Percentual de arestas existentes em relação ao grafo completo $K_{|V|}$, variando no intervalo $[0, 1]$:
    \begin{equation}
        D = \frac{2 \cdot |E|}{|V| \cdot (|V| - 1)}
    \end{equation}
    \item \textbf{Vizinhos e Grau de um Vértice:} O grau $d(v)$ corresponde à cardinalidade do conjunto de adjacências de $v$, obtido varrendo a linha $v$ da matriz.
\end{itemize}

% ==============================================================================
% 4. CONECTIVIDADE E BUSCA (MÓDULO DE PEDRO)
% ==============================================================================
\section{Módulo de Conectividade e Busca}

Esta seção descreve as quatro funcionalidades que constituem o núcleo de conectividade da biblioteca.

\subsection{Busca em Largura (BFS), Árvore de Busca e Arestas Fora da Árvore}
A Busca em Largura (Breadth-First Search) percorre a rede em ondas concêntricas a partir de um vértice de partida $s \in V$. A exploração é gerenciada por uma fila FIFO (\texttt{collections.deque}), garantindo complexidade temporal $O(V + E)$.

Durante a travessia na componente conexa de $s$, as arestas encontradas são estritamente particionadas em duas categorias:
\begin{enumerate}
    \item \textbf{Arestas da Árvore de Busca ($E_T$):} Arestas $(u, v)$ utilizadas para descobrir e alcançar pela primeira vez um vértice $v$ ainda não visitado a partir de $u$. Essas arestas formam a árvore geradora da componente.
    \item \textbf{Arestas Fora da Árvore ($E_{\text{fora}}$):} Arestas existentes entre pares de vértices que já foram alcançados pela busca, mas que não foram utilizadas como caminho de descoberta. Formalmente conhecidas como \textit{arestas de cruzamento} (\textit{cross-edges}), essas arestas são as responsáveis diretas pelo fechamento de ciclos na rede.
\end{enumerate}

\begin{lstlisting}[caption={Implementação da Busca em Largura em \texttt{conectividade.py}}, label={cod:bfs}]
def busca_em_largura(grafo, vertice_inicial):
    visitados = [vertice_inicial]
    visitados_set = {vertice_inicial}
    fila = deque([vertice_inicial])

    arvore = []
    arvore_arestas = set()

    while fila:
        u = fila.popleft()
        for viz in grafo.vizinhos(u):
            if viz not in visitados_set:
                visitados_set.add(viz)
                visitados.append(viz)
                fila.append(viz)

                arvore.append((u, viz))
                arvore_arestas.add((min(u, viz), max(u, viz)))

    # Coleta arestas entre visitados que nao pertencem a arvore
    fora_arvore = []
    for u in visitados:
        for viz in grafo.vizinhos(u):
            if u < viz and viz in visitados_set:
                aresta = (u, viz)
                if aresta not in arvore_arestas:
                    fora_arvore.append(aresta)

    return visitados, arvore, fora_arvore
\end{lstlisting}

\subsection{Componentes Conexas via Algoritmo de Roy (Fechamento Transitivo)}
Conforme sugestão do docente, o cálculo das componentes conexas foi implementado através do \textbf{Algoritmo de Bernard Roy (1959)} (Fechamento Transitivo).

A relação de conectividade em um grafo não direcionado constitui uma relação de equivalência (reflexiva, simétrica e transitiva). O algoritmo calcula a matriz booleana de alcançabilidade $R$, onde $R_{ij} = \text{True}$ se e somente se existe uma cadeia ligando o vértice $i$ ao vértice $j$. A recorrência algébrica de Roy é dada por:
\begin{equation}
    R^{(k)}_{ij} = R^{(k-1)}_{ij} \lor \left( R^{(k-1)}_{ik} \land R^{(k-1)}_{kj} \right), \quad \forall k \in \{1, \dots, N\}
\end{equation}

Após a convergência das $N$ iterações, os vértices cujas linhas possuem valores coincidentes pertencem à mesma classe de equivalência (componente conexa). A listagem \ref{cod:roy} apresenta a implementação direta e intuitiva do método.

\begin{lstlisting}[caption={Algoritmo de Roy para Componentes Conexas}, label={cod:roy}]
def componentes_conexas_roy(grafo):
    n = grafo.ordem()
    if n == 0: return 0, []

    # 1. Matriz de alcancabilidade R
    R = [[False] * (n + 1) for _ in range(n + 1)]
    for i in range(1, n + 1):
        R[i][i] = True
        for viz in grafo.vizinhos(i):
            R[i][viz] = True

    # 2. Fechamento Transitivo de Roy (3 loops tradicionais)
    for k in range(1, n + 1):
        for i in range(1, n + 1):
            if R[i][k]:
                for j in range(1, n + 1):
                    if R[k][j]:
                        R[i][j] = True

    # 3. Agrupamento em componentes
    visitados = set()
    componentes = []
    for i in range(1, n + 1):
        if i not in visitados:
            comp = [j for j in range(1, n + 1) if R[i][j]]
            visitados.update(comp)
            componentes.append(comp)

    return len(componentes), componentes
\end{lstlisting}

\subsection{Identificação de Vértices de Articulação (Pontos de Corte)}
Um vértice $v \in V$ é denominado \textbf{vértice de articulação} se a sua supressão (juntamente com as arestas a ele incidentes) desconecta o subgrafo $G - \{v\}$, aumentando estritamente o número de componentes conexas da rede.

A verificação foi implementada com alta eficiência explorando propriedades estruturais:
\begin{itemize}
    \item Vértices de grau 0 ou grau 1 nunca são pontos de articulação (remover folhas não desconecta os nós remanescentes).
    \item Para um vértice $v$ com vizinhos $Adj(v) = \{w_1, w_2, \dots, w_d\}$ com $d \ge 2$: executa-se uma BFS no subgrafo $G - \{v\}$ partindo de $w_1$. Se a busca for incapaz de alcançar qualquer um dos outros vizinhos $w_i$, conclui-se que $v$ era a ponte obrigatória de passagem entre eles. Portanto, $v$ é articulação.
\end{itemize}

\begin{lstlisting}[caption={Verificação de Vértice de Articulação}, label={cod:articulacao}]
def verificar_articulacao(grafo, vertice):
    vizinhos = grafo.vizinhos(vertice)
    if len(vizinhos) <= 1:
        return False

    # BFS no subgrafo excluindo o vertice testado
    visitados = {vizinhos[0]}
    fila = deque([vizinhos[0]])

    while fila:
        u = fila.popleft()
        for viz in grafo.vizinhos(u):
            if viz != vertice and viz not in visitados:
                visitados.add(viz)
                fila.append(viz)

    # Se algum vizinho nao foi alcancado, o vertice eh articulacao
    for viz in vizinhos[1:]:
        if viz not in visitados:
            return True

    return False
\end{lstlisting}

\subsection{Detecção de Ciclos}
A presença de ciclos indica a existência de rotas redundantes ou circuitos fechados entre usuários. A função \texttt{possui\_ciclo} percorre o grafo via BFS mantendo o rastreamento explícito do predecessor imediato $(\text{atual}, \text{pai})$. Caso um vizinho já visitado seja encontrado e ele seja distinto do nó pai, um caminho alternativo foi detectado, comprovando a existência de ciclo e retornando \texttt{True} antecipadamente.

% ==============================================================================
% 5. CAMINHOS MÍNIMOS (DIJKSTRA)
% ==============================================================================
\section{Módulo de Caminhos Mínimos}

O cálculo de distâncias mínimas foi baseado no \textbf{Algoritmo de Dijkstra}, utilizando fila de prioridade implementada via \textit{Min-Heap} binário (\texttt{heapq}).

Como todos os pesos das arestas na rede social representam custos reais positivos ($w(e) > 0$), a propriedade de subestrutura ótima é garantida. O algoritmo mantém um vetor de menores distâncias acumuladas $dist[u]$ e um mapa de predecessores $pred[u]$, permitindo recuperar não apenas o custo numérico mínimo, mas a \textbf{rota completa passo a passo} entre o vértice de origem e todos os outros vértices da rede.

% ==============================================================================
% 6. RESULTADOS EXPERIMENTAIS E VALIDAÇÃO DAS INSTÂNCIAS
% ==============================================================================
\section{Resultados Experimentais e Validação das Instâncias}

Para validação completa da biblioteca, foram avaliadas as quatro instâncias de teste exigidas nas especificações do trabalho.

\subsection{Instância 1 -- Grafo Pequeno e Conexo ($|V|=10, |E|=13$)}
A Instância 1 contempla um grafo denso e conexo de pequeno porte, projetado para validação manual dos resultados.

% PLACEHOLDER DE IMAGEM: GRAFO DA INSTANCIA 1
\begin{figure}[H]
    \centering
    % Substitua o caminho abaixo pela imagem gerada do grafo (ex: imagens/grafo_instancia1.png)
    \fbox{\parbox[c][5cm][c]{0.8\textwidth}{\centering \textbf{[Inserir Aqui a Imagem/Diagrama do Grafo da Instância 1]}\\ \small Sugestão de arquivo: \texttt{imagens/instancia1.png}}}
    \caption{Topologia da rede social para a Instância 1 ($|V|=10, |E|=13$).}
    \label{fig:instancia1}
\end{figure}

\begin{table}[H]
    \centering
    \caption{Conferência dos Resultados Obtidos na Instância 1}
    \begin{tabular}{@{}lll@{}}
        \toprule
        \textbf{Funcionalidade} & \textbf{Resultado Obtido no Software} & \textbf{Conferência Analítica} \\ \midrule
        Ordem ($|V|$) & 10 & 10 vértices cadastrados \\
        Tamanho ($|E|$) & 13 & 13 conexões únicas \\
        Densidade $\varepsilon(G)$ & 1.3000 & $13 / 10 = 1.3000$ \\
        Densidade Relativa & 0.2889 & $2 \times 13 / (10 \times 9) = 0.2889$ \\
        Componentes Conexas & 1 componente com 10 vértices & Grafo totalmente conexo \\
        Vértices de Articulação & Nenhum & Grafo biconexo (2-conexo) \\
        Detecção de Ciclos & Possui Ciclo (True) & $|E| = 13 > 10 - 1 = 9$ (Obrigatório) \\
        BFS a partir de 1 & 10 visitados; 9 na árvore; 4 fora & $9 + 4 = 13$ arestas totais \\ \bottomrule
    \end{tabular}
    \label{tab:res_instancia1}
\end{table}

A sequência de visitação da BFS iniciada no vértice 1 produziu a árvore geradora com as seguintes 9 arestas:
$\{(1,2), (1,3), (2,4), (2,5), (3,6), (4,7), (5,8), (7,9), (8,10)\}$.
As 4 arestas que fecham ciclos e ficaram fora da árvore de busca foram identificadas como:
$\{(3,5), (5,7), (6,8), (9,10)\}$.

\subsection{Instância 2 -- Grafo Desconexo ($|V|=12, |E|=12$)}
A Instância 2 avalia a robustez dos módulos na presença de múltiplas componentes conexas e pontos críticos de vulnerabilidade.

% PLACEHOLDER DE IMAGEM: GRAFO DA INSTANCIA 2
\begin{figure}[H]
    \centering
    % Substitua o caminho abaixo pela imagem do grafo da instancia 2
    \fbox{\parbox[c][5cm][c]{0.8\textwidth}{\centering \textbf{[Inserir Aqui a Imagem/Diagrama do Grafo da Instância 2]}\\ \small Sugestão de arquivo: \texttt{imagens/instancia2.png}}}
    \caption{Grafo desconexo da Instância 2 dividido em 3 comunidades isoladas.}
    \label{fig:instancia2}
\end{figure}

O algoritmo de Roy identificou com exatidão as 3 componentes conexas:
\begin{itemize}
    \item \textbf{Componente 1:} $\{1, 2, 3, 4, 5\}$ ($|E|=5$, contém ciclo 1-2-3 e caminho 2-4-5).
    \item \textbf{Componente 2:} $\{6, 7, 8, 9\}$ ($|E|=4$, ciclo $C_4$).
    \item \textbf{Componente 3:} $\{10, 11, 12\}$ ($|E|=3$, ciclo $C_3$).
\end{itemize}

Na análise de articulação, o sistema apontou com precisão os vértices \textbf{2} e \textbf{4}:
\begin{itemize}
    \item A remoção do vértice 2 fragmenta a Componente 1 em dois subconjuntos isolados: $\{1, 3\}$ e $\{4, 5\}$.
    \item A remoção do vértice 4 isola o usuário 5 do restante da rede.
\end{itemize}

A BFS iniciada no vértice 1 limitou-se rigorosamente à Componente 1, gerando 4 arestas de árvore e identificando a aresta $(2, 3)$ como a única fora da árvore.

\subsection{Instâncias 3 e 4 -- Avaliação de Escalabilidade e Desempenho}
Para avaliar a eficiência da biblioteca frente ao aumento de escala, foram processadas as instâncias 3 (grafo esparso) e 4 (grafo denso de grande porte).

\begin{table}[H]
    \centering
    \caption{Métricas e Desempenho nas Instâncias de Maior Escala}
    \begin{tabular}{@{}lcccc@{}}
        \toprule
        \textbf{Instância} & \textbf{Ordem ($|V|$)} & \textbf{Tamanho ($|E|$)} & \textbf{Densidade $\varepsilon(G)$} & \textbf{Tempo Roy / BFS} \\ \midrule
        Instância 3 (Esparso) & 50 & 53 & 1.0600 & $< 1.5\text{ ms}$ \\
        Instância 4 (Grande/Denso) & 500 & 6.717 & 13.4340 & $\approx 65\text{ ms}$ \\ \bottomrule
    \end{tabular}
    \label{tab:desempenho}
\end{table}

\textbf{Discussão de Desempenho:} A utilização da matriz de adjacência impõe consumo espacial de ordem $O(N^2)$. Para a Instância 4 com $N=500$, a alocação em memória é de $501 \times 501 \approx 251.000$ ponteiros, perfeitamente comportada pela memória primária. O tempo de resposta das buscas e checagens permaneceu na ordem de milissegundos, demonstrando a adequação da implementação aos limites estipulados.

% ==============================================================================
% 7. GUIA DE EXECUÇÃO DO SOFTWARE
% ==============================================================================
\section{Instruções de Compilação e Execução}

O projeto foi construído utilizando exclusivamente a biblioteca padrão do Python 3 (sem necessidade de instalação de dependências via \texttt{pip}).

Para executar o programa interativo principal:
\begin{lstlisting}[language=bash, numbers=none]
$ python main.py
\end{lstlisting}

O usuário será solicitado a informar o caminho do arquivo de teste (exemplo: \texttt{instancias/instancia1.txt}) e poderá navegar pelas opções de 1 a 10 no menu de console.

Para executar a suíte independente de testes automatizados do módulo de conectividade:
\begin{lstlisting}[language=bash, numbers=none]
$ python tests_independentes/test_pedro.py
\end{lstlisting}

% ==============================================================================
% 8. CONCLUSÃO
% ==============================================================================
\section{Conclusão}

O desenvolvimento do Trabalho Prático I permitiu consolidar a aplicação dos fundamentos de Teoria dos Grafos na modelagem de sistemas computacionais reais. A representação por matriz de adjacência ponderada mostrou-se eficaz na rápida consulta de adjacências e no cálculo algébrico do fechamento transitivo de Roy.

A modularização estrita do código permitiu que os três discentes atuassem de maneira integrada e desacoplada, garantindo alta legibilidade e robustez nas validações empíricas e teóricas das instâncias avaliadas.

% ==============================================================================
% REFERÊNCIAS BIBLIOGRÁFICAS
% ==============================================================================
\begin{thebibliography}{99}
    \bibitem{cormen} CORMEN, T. H.; LEISERSON, C. E.; RIVEST, R. L.; STEIN, C. \textbf{Algoritmos: Teoria e Prática}. 3. ed. Rio de Janeiro: Elsevier, 2012.
    \bibitem{szwarcfiter} SZWARCFITER, J. L. \textbf{Grafos e Algoritmos Computacionais}. Rio de Janeiro: Campus, 1984.
    \bibitem{bondy} BONDY, J. A.; MURTY, U. S. R. \textbf{Graph Theory with Applications}. London: Macmillan, 1976.
    \bibitem{roy} ROY, B. \textbf{Transitivité et connexité}. Comptes Rendus de l'Académie des Sciences de Paris, v. 249, p. 216--218, 1959.
\end{thebibliography}

\end{document}
